%%=============================================================================
%% Voorwoord
%%=============================================================================

\chapter*{\IfLanguageName{dutch}{Woord vooraf}{Preface}}%
\label{ch:voorwoord}

Met het schrijven van dit voorwoord leg ik de laatste hand aan mijn bachelorproef, het sluitstuk van mijn opleiding Toegepaste Informatica aan HoGent. 

Toen ik aan het traject van AI \& Data begon, wist ik zeker: ik wilde een onderwerp uitdiepen waarin data centraal stond. Het idee om orde te scheppen in de chaos van complexe en gefragmenteerde bedrijfsdata heeft me altijd gefascineerd. De uitdaging om voor IPCOM een Proof of Concept te bouwen voor een in-house Master Data Management-oplossing binnen Microsoft Fabric, bleek na een tijdje een onderwerp voor mij waar ik echt interesse in had. Het bleek een boeiende technische puzzel waarbij ik niet alleen diep in PySpark, data matching en lakehouse-architecturen mocht duiken, maar ook leerde hoe deze theorieën zich vertalen naar de realiteit van een grote onderneming.

Natuurlijk is deze bachelorproef niet uitsluitend het resultaat van mijn eigen inspanningen. Ik wil dan ook graag de tijd nemen om een aantal mensen te bedanken die een cruciale rol hebben gespeeld in de totstandkoming van dit werk.

In de eerste plaats wil ik mijn stagebedrijf, IPCOM, bedanken voor de mogelijkheid tot het uitbouwen van deze oplossing en het vertrouwen. Een speciale dankjewel aan mijn co-promotor, Gert-Jan Dugardein, voor de expertise en begeleiding binnen een voor mij eerst onbekende wereld en de verhelderende brainstormsessies. Bedankt voor het geduld en de inzichten, zelfs wanneer ik een bepaald concept niet direct snapte en het paar keer opnieuw moest vragen.

Daarnaast wil ik mijn academische promotor, Karine Van Driessche, oprecht bedanken. Uw kritische blik voor het schrijven van een bachelorproef, de constructieve feedback en de gerichte bijsturingen hebben ervoor gezorgd dat dit onderzoek vooral academisch op punt staat.

Ik wens u veel leesplezier.

\vspace{2cm}

\noindent Alvaro Van Damme\\
\noindent Wetteren, \today